% =====================================================================
%  MANUAL TÉCNICO — ATARA
%  Sistema de Alerta Temprana y Atención al Rendimiento Académico
%  División de Educación Básica — Ministerio de Educación Pública (MEP)
%  Curso: EIF408 — Proyecto de Ingeniería de Software
%  Fecha: Abril 2026
% =====================================================================
\documentclass[11pt,a4paper,oneside]{article}

% --------------------- Codificación e idioma ---------------------------
\usepackage[utf8]{inputenc}
\usepackage[T1]{fontenc}
\usepackage[spanish,es-tabla]{babel}
\usepackage{lmodern}
\usepackage{microtype}

% --------------------- Geometría y estilo ------------------------------
\usepackage[a4paper,margin=2.5cm,headheight=14pt]{geometry}
\usepackage{fancyhdr}
\usepackage{titlesec}
\usepackage{parskip}

% --------------------- Tablas y listas ---------------------------------
\usepackage{booktabs}
\usepackage{longtable}
\usepackage{array}
\usepackage{tabularx}
\usepackage{multirow}
\usepackage{enumitem}

% --------------------- Color, cajas, código ----------------------------
\usepackage{xcolor}
\usepackage[most]{tcolorbox}
\usepackage{listings}

% --------------------- URLs y rutas de API (quiebre en monospace) -----
\usepackage[hyphens]{url}
\def\UrlBreaks{\do\/\do\{\do\}\do\_\do\-\do\=\do\~\do\.}

% --------------------- Diagramas ---------------------------------------
\usepackage{tikz}
\usetikzlibrary{shapes.geometric, arrows.meta, positioning, fit, backgrounds, calc}

% --------------------- Hipervínculos -----------------------------------
\usepackage[hidelinks,colorlinks=true,linkcolor=ataraAzul,urlcolor=ataraAzul,citecolor=ataraAzul]{hyperref}

% --------------------- Paleta de colores ATARA -------------------------
\definecolor{ataraAzul}{HTML}{1F3A93}
\definecolor{ataraAzulClaro}{HTML}{4A6FBF}
\definecolor{ataraGris}{HTML}{F4F6F8}
\definecolor{ataraGrisOscuro}{HTML}{4A4A4A}
\definecolor{ataraVerde}{HTML}{1B7F3B}
\definecolor{ataraNaranja}{HTML}{D97E1F}
\definecolor{ataraRojo}{HTML}{B5292F}
\definecolor{ataraFondoCodigo}{HTML}{F8F8F8}

% --------------------- Estilo de títulos -------------------------------
\titleformat{\section}
  {\color{ataraAzul}\Large\bfseries}{\thesection.}{0.6em}{}[\vspace{-0.4em}\color{ataraAzul}\rule{\linewidth}{0.6pt}]
\titleformat{\subsection}
  {\color{ataraAzul}\large\bfseries}{\thesubsection.}{0.5em}{}
\titleformat{\subsubsection}
  {\color{ataraAzulClaro}\normalsize\bfseries}{\thesubsubsection.}{0.5em}{}

% --------------------- Encabezado y pie --------------------------------
\pagestyle{fancy}
\fancyhf{}
\fancyhead[L]{\textbf{\color{ataraAzul}ATARA} \textcolor{ataraGrisOscuro}{— Manual Técnico}}
\fancyhead[R]{\textcolor{ataraGrisOscuro}{EIF408 \textperiodcentered{} 2026}}
\fancyfoot[C]{\thepage}
\renewcommand{\headrulewidth}{0.4pt}
\renewcommand{\headrule}{\hbox to\headwidth{\color{ataraAzul}\leaders\hrule height \headrulewidth\hfill}}

% --------------------- Configuración de listings -----------------------
\lstset{
  basicstyle=\ttfamily\footnotesize,
  backgroundcolor=\color{ataraFondoCodigo},
  frame=single,
  rulecolor=\color{ataraAzulClaro!40},
  framesep=4pt,
  xleftmargin=8pt,
  xrightmargin=4pt,
  numbers=none,
  showstringspaces=false,
  breaklines=true,
  breakatwhitespace=false,
  keywordstyle=\color{ataraAzul}\bfseries,
  commentstyle=\color{ataraGrisOscuro}\itshape,
  stringstyle=\color{ataraVerde},
  captionpos=b,
  literate=
    {á}{{\'a}}1 {é}{{\'e}}1 {í}{{\'\i}}1 {ó}{{\'o}}1 {ú}{{\'u}}1
    {Á}{{\'A}}1 {É}{{\'E}}1 {Í}{{\'I}}1 {Ó}{{\'O}}1 {Ú}{{\'U}}1
    {ñ}{{\~n}}1 {Ñ}{{\~N}}1 {ü}{{\"u}}1 {Ü}{{\"U}}1 {¿}{{?`}}1 {¡}{{!`}}1
}

\lstdefinestyle{json}{language=Java,morekeywords={true,false,null}}
\lstdefinestyle{sql}{language=SQL}
\lstdefinestyle{java}{language=Java}
\lstdefinestyle{bash}{language=bash,morekeywords={mvnw,docker,docker-compose,curl}}

% --------------------- Cajas tcolorbox ---------------------------------
\newtcolorbox{ataracaja}[1][]{%
  colback=ataraAzul!4,
  colframe=ataraAzul!50,
  arc=2mm,
  boxrule=0.6pt,
  fonttitle=\bfseries,
  coltitle=ataraAzul,
  #1
}

\newtcolorbox{ataraInfo}{
  colback=ataraAzul!4, colframe=ataraAzul!40,
  arc=2mm, boxrule=0.6pt
}
\newtcolorbox{ataraAlerta}{
  colback=ataraRojo!5, colframe=ataraRojo!60,
  arc=2mm, boxrule=0.6pt
}
\newtcolorbox{ataraNota}{
  colback=ataraNaranja!8, colframe=ataraNaranja!60,
  arc=2mm, boxrule=0.6pt
}
\newtcolorbox{ataraExito}{
  colback=ataraVerde!7, colframe=ataraVerde!60,
  arc=2mm, boxrule=0.6pt
}

% --------------------- Comandos utilitarios ----------------------------
\newcommand{\code}[1]{\texttt{\small #1}}
% \ruta usa \path (paquete url) -> quiebra lineas en / { } _ = sin escapar
\newcommand{\ruta}[1]{\small\path{#1}}
\newcommand{\pass}{\textbf{\color{ataraVerde}PASS}}
\newcommand{\fail}{\textbf{\color{ataraRojo}FAIL}}

% =====================================================================
%                                 PORTADA
% =====================================================================
\begin{document}
\pagenumbering{roman}
\thispagestyle{empty}

\begin{center}
\vspace*{1.5cm}

{\Huge\bfseries\color{ataraAzul} ATARA}\\[0.5em]
{\Large Sistema de Alerta Temprana y Atención al Rendimiento Académico}\\[0.4em]
{\normalsize\itshape División de Educación Básica — MEP, Costa Rica}

\vspace{2.2cm}
\rule{0.85\linewidth}{1.5pt}

\vspace{0.7cm}
{\Huge\bfseries Manual Técnico}\\[0.4em]
{\Large de la Plataforma Completa}\\[0.4em]
{\large\itshape Backend REST API + Frontend Vite}

\vspace{0.5cm}
\rule{0.85\linewidth}{1.5pt}

\vspace{2.5cm}

\begin{tabular}{rl}
\textbf{Curso:}              & EIF408 — Proyecto de Ingeniería de Software \\[0.2em]
\textbf{Versión del API:}    & 0.0.1-SNAPSHOT \\[0.2em]
\textbf{Fecha de emisión:}   & Abril 2026 \\[0.2em]
\textbf{Stack principal:}    & Spring Boot 4.0.3 \textperiodcentered{} Java 17 \textperiodcentered{} PostgreSQL 17 \\[0.2em]
\textbf{Stack frontend:}     & Vite 5 \textperiodcentered{} JavaScript Vanilla \textperiodcentered{} Chart.js 4 \\[0.2em]
\textbf{Puerto API local:}   & \code{localhost:8081} \\[0.2em]
\textbf{Puerto frontend:}    & \code{localhost:3000} \\[0.2em]
\textbf{Repositorio:}        & ATARA-Backend (rama \code{main}) \\
\end{tabular}

\vspace{1.5cm}

\begin{ataracaja}[title=Alcance del documento]
Este manual cubre la \textbf{plataforma completa} de ATARA hasta abril 2026: arquitectura,
modelo de datos, los \textbf{14 controladores REST}, los \textbf{7 módulos funcionales},
el frontend Vite, la infraestructura de pruebas funcionales y la configuración de despliegue
en producción. Reemplaza al borrador previo y se mantiene como documento técnico de referencia.
\end{ataracaja}

\end{center}
\newpage

% =====================================================================
%                                 ÍNDICE
% =====================================================================
\tableofcontents
\newpage

\pagenumbering{arabic}

% =====================================================================
\section{Introducción}
% =====================================================================

\subsection{Propósito del sistema}

ATARA (\emph{Sistema de Alerta Temprana y Atención al Rendimiento Académico}) es una plataforma
web diseñada para la \textbf{División de Educación Básica del Ministerio de Educación Pública
de Costa Rica} (MEP). Su objetivo es detectar, de manera temprana y automatizada, dificultades
académicas en estudiantes de primaria mediante un modelo de evaluación por saberes
(Conceptual, Procedimental, Actitudinal) y la generación de \textbf{alertas temáticas} cuando
el desempeño promedio de un estudiante en un eje temático cae por debajo del umbral esperado.

\subsection{Audiencia}

\begin{itemize}[leftmargin=*]
  \item \textbf{Equipo de desarrollo y mantenimiento} — para entender la arquitectura,
        integrar nuevas funcionalidades y depurar incidencias.
  \item \textbf{Equipo de despliegue / DevOps} — para configurar entornos, gestionar
        variables de entorno, ejecutar migraciones y monitorear el sistema.
  \item \textbf{Coordinadores académicos del MEP} — como referencia formal del modelo
        funcional implementado.
  \item \textbf{Auditores y revisores docentes} — para validar el cumplimiento del modelo
        pedagógico y del control de acceso.
\end{itemize}

\subsection{Convenciones del documento}

\begin{itemize}[leftmargin=*]
  \item \code{Código en monoespaciado} indica nombres de clases, archivos, endpoints, variables
        de entorno, identificadores de configuración o cuerpos JSON.
  \item Los \textbf{verbos HTTP} (GET, POST, PUT, DELETE) se escriben en mayúscula seguidos de
        la ruta.
  \item Las cajas con borde azul son \emph{informativas}; con borde naranja son \emph{notas
        de implementación}; con borde rojo son \emph{advertencias}; con borde verde son
        \emph{evidencias o confirmaciones}.
\end{itemize}

% =====================================================================
\section{Visión general del sistema}
% =====================================================================

\subsection{Stack tecnológico}

\begin{table}[h]
\centering
\renewcommand{\arraystretch}{1.2}
\begin{tabular}{@{}p{4.0cm} p{6.0cm} p{4.0cm}@{}}
\toprule
\textbf{Componente} & \textbf{Tecnología} & \textbf{Versión} \\
\midrule
Framework backend       & Spring Boot                    & 4.0.3 \\
Lenguaje backend        & Java                           & 17 (LTS) \\
Base de datos           & PostgreSQL                     & 17 (Docker) / 15-alpine (tests) \\
Migraciones de esquema  & Flyway (PostgreSQL)            & integrado con Spring Boot \\
Seguridad               & Spring Security + JJWT         & 0.12.6 \\
Validación              & Jakarta Bean Validation        & integrada \\
Persistencia            & Spring Data JPA + Hibernate    & integrada \\
OCR                     & Apache PDFBox + Tess4j         & 3.0.3 / 5.11.0 \\
Pruebas funcionales     & JUnit 5 + Testcontainers       & 1.20.4 \\
Aserciones              & AssertJ                        & integrada \\
Cliente HTTP de pruebas & Spring \code{RestTemplate}     & integrada \\
Build tool              & Apache Maven Wrapper           & \code{mvnw} \\
Frontend                & Vite + Vanilla JS              & Vite 5 \\
Gráficas frontend       & Chart.js                       & 4.5.1 \\
Contenedor de BD        & Docker Compose                 & \code{docker-compose.yml} \\
\bottomrule
\end{tabular}
\caption{Componentes del stack tecnológico oficial de ATARA.}
\end{table}

\subsection{Arquitectura en capas}

El backend sigue una arquitectura clásica en capas con separación estricta de
responsabilidades. Las flechas representan la dirección permitida de las dependencias.

\begin{center}
\begin{tikzpicture}[
  node distance=0.6cm and 1.0cm,
  capa/.style={rectangle, rounded corners=2pt, draw=ataraAzul!60, fill=ataraAzul!8,
               minimum width=3.5cm, minimum height=1.05cm, align=center, font=\small\bfseries},
  flecha/.style={-{Stealth[length=2.5mm]}, thick, ataraAzul!70},
  rol/.style={font=\scriptsize\itshape, ataraGrisOscuro, align=center}
]
\node[capa] (cli) {Cliente / Frontend\\\scriptsize Vite + Vanilla JS};
\node[capa, right=of cli] (ctrl) {Controller\\\scriptsize @Valid, DTO};
\node[capa, right=of ctrl] (svc) {Service\\\scriptsize Reglas de negocio};
\node[capa, right=of svc] (repo) {Repository\\\scriptsize Spring Data JPA};
\node[capa, right=of repo] (db) {Modelo / BD\\\scriptsize Hibernate $\to$ PostgreSQL};

\draw[flecha] (cli) -- (ctrl) node[midway, above, rol] {HTTPS\\JSON};
\draw[flecha] (ctrl) -- (svc) node[midway, above, rol] {DTO};
\draw[flecha] (svc) -- (repo) node[midway, above, rol] {Entidades};
\draw[flecha] (repo) -- (db) node[midway, above, rol] {SQL};
\end{tikzpicture}
\end{center}

\textbf{Responsabilidades por capa:}

\begin{itemize}[leftmargin=*]
  \item \textbf{Controller} — recibe peticiones HTTP, aplica \code{@Valid} sobre los DTO
        de entrada, delega al servicio y responde con \code{ResponseEntity<ApiResponse<T>>}.
        \emph{Nunca} contiene lógica de negocio ni accede a repositorios.
  \item \textbf{Service} — implementa toda la lógica de negocio, validaciones de reglas,
        transformación entidad$\leftrightarrow$DTO y verificaciones de control de acceso
        a través del \code{ContextoUsuario}.
  \item \textbf{Repository} — extiende \code{JpaRepository<Entity, Long>} (Spring Data JPA).
        Solo realiza consultas a la base de datos.
  \item \textbf{Modelo (entidades JPA)} — todas las relaciones se declaran con
        \code{FetchType.LAZY} para evitar materialización de grafos innecesarios.
  \item \textbf{Seguridad} — el filtro \code{JwtAuthenticationFilter} valida el token en cada
        request; el \code{ContextoUsuarioService} construye un \code{ContextoUsuario} inmutable
        para cada request a partir del usuario autenticado.
\end{itemize}

\subsection{Estructura de paquetes Java}

\begin{lstlisting}[style=bash,caption={Distribución del código Java}]
com.atara.deb.ataraapi
|-- controller/   14 controladores REST
|-- service/      Interfaces de servicio
|   `-- impl/     Implementaciones (Impl)
|-- repository/   Repositorios Spring Data JPA
|-- model/        27 entidades JPA
|   `-- enums/    Enumeraciones del dominio
|-- dto/          DTOs por dominio
|   |-- alerta/   |-- estudiante/  |-- piad/
|   |-- auth/     |-- evaluacion/  |-- reporte/
|   |-- saber/    |-- matricula/   |-- usuario/
|   `-- visualizacion/
|-- security/     SecurityConfig, JwtService,
|                 JwtAuthenticationFilter,
|                 ContextoUsuario, ContextoUsuarioService
|-- exception/    GlobalExceptionHandler,
|                 AccesoDenegadoException,
|                 TokenRefreshException
`-- AtaraApiApplication.java
\end{lstlisting}

\subsection{Estructura del frontend}

El frontend (\code{temp-frontend/}) es una aplicación Vite + JavaScript vanilla, sin framework
SPA. Está pensada para integrarse con el backend mediante el cliente \code{api.js}, que
auto-extrae el campo \code{.data} del envoltorio \code{ApiResponse<T>}.

\begin{lstlisting}[style=bash,caption={Estructura del frontend Vite}]
temp-frontend/
|-- index.html              entrada principal (carga de pagina dinamica)
|-- vite.config.js          puerto 3000, proxy /api -> :8081
|-- package.json            dependencias (Chart.js)
`-- src/
    |-- main.js             ruteo por query params
    |-- api.js              cliente HTTP, manejo de JWT, unwrap ApiResponse
    |-- confirm.js          dialogos de confirmacion personalizados
    |-- pedagogicaData.js   constantes pedagogicas estaticas
    |-- style.css           estilos globales
    `-- pages/
        |-- aniosLectivos.js     |-- alertasTempranas.js
        |-- estudiantes.js       |-- evaluacionesSaber.js
        |-- matriculas.js        |-- visualizaciones.js
        |-- evaluaciones.js      |-- reportes.js
        |-- alertas.js           |-- importarPiad.js
        |-- secciones.js         |-- sesion.js
\end{lstlisting}

% =====================================================================
\section{Modelo de datos y migraciones}
% =====================================================================

\subsection{Gobierno del esquema}

El esquema de base de datos es \textbf{propiedad exclusiva de Flyway}. Hibernate opera con
\code{spring.jpa.hibernate.ddl-auto=none}: nunca se permite que la JPA modifique la estructura
de la BD. La única vía válida para cambiar el esquema es crear una nueva migración versionada
(\code{V12\_\_descripcion.sql}, \code{V13\_\_...}, etc.).

\begin{ataraAlerta}
\textbf{Restricción crítica:} las migraciones existentes son \emph{inmutables}. Modificarlas
hace que Flyway falle con un \code{Validate failed: checksum mismatch} y la aplicación no arranca.
Toda corrección de esquema o de datos debe ir en una migración nueva.
\end{ataraAlerta}

\subsection{Cronología de migraciones}

\begin{table}[h]
\centering
\renewcommand{\arraystretch}{1.15}
\begin{tabular}{@{}p{1.0cm} p{12.5cm}@{}}
\toprule
\textbf{Versión} & \textbf{Contenido} \\
\midrule
V1  & Esquema inicial: 20 tablas base, más de 40 índices, 13 \emph{triggers} de auditoría
       (\code{registro\_auditoria} con cuerpo JSONB) y la función
       \code{fn\_actualizar\_updated\_at}. \\
V2  & Datos semilla: usuarios, centros educativos, niveles, secciones, estudiantes,
       matrículas y configuración pedagógica básica. \\
V3  & 8 vistas de consulta para reportería:
       \code{vw\_criterios\_completos}, \code{vw\_rendimiento\_periodo\_activo},
       entre otras. \\
V4  & Módulo de saberes: tablas \code{tipos\_saber}, \code{ejes\_tematicos},
       \code{niveles\_desempeno}, \code{evaluaciones\_saber},
       \code{detalle\_evaluacion\_saber}, \code{alertas\_tematicas} y dos vistas. \\
V5  & Estudiantes adicionales para enriquecer el set de pruebas. \\
V6  & Soporte multi-materia: agrega \code{materia\_id} a \code{ejes\_tematicos},
       \code{evaluaciones\_saber} y \code{alertas\_tematicas}; añade 63 ejes nuevos
       (Matemáticas, Ciencias, Estudios Sociales) y actualiza las vistas. \\
V7  & Restricción única en \code{alertas\_tematicas} por
       \code{(estudiante\_id, periodo\_id, eje\_tematico\_id, materia\_id)}. \\
V8  & Tabla \code{usuario\_materias}: relación N:M entre docentes y materias asignadas. \\
V9  & Inserta el docente de prueba \code{kcortes@atara.mep.go.cr} (Keylor Cortés). \\
V10 & Asignación de materias adicionales (Español y Ciencias) para el docente
       \code{avargas@atara.mep.go.cr}. \\
V11 & Corrección de la asignación del docente de la sección 1°D. \\
\bottomrule
\end{tabular}
\caption{Cronología de migraciones Flyway aplicadas (V1–V11).}
\end{table}

\subsection{Diagrama lógico del modelo}

\begin{center}
\begin{tikzpicture}[
  ent/.style={rectangle, rounded corners=1.5pt, draw=ataraAzul!50,
              fill=ataraAzul!6, font=\scriptsize\bfseries,
              minimum height=0.6cm, minimum width=2.0cm, align=center},
  rel/.style={-{Stealth[length=1.6mm]}, thin, ataraAzul!60},
  node distance=0.55cm and 0.7cm
]
\node[ent] (cu) {Centro\\Educativo};
\node[ent, right=of cu] (al) {AnioLectivo};
\node[ent, right=of al] (per) {Periodo};
\node[ent, right=of per] (mat) {Materia};

\node[ent, below=of cu] (sec) {Seccion};
\node[ent, right=of sec] (us) {Usuario};
\node[ent, right=of us] (ts) {TipoSaber};
\node[ent, right=of ts] (et) {EjeTematico};

\node[ent, below=of sec] (est) {Estudiante};
\node[ent, right=of est] (mtr) {Matricula};
\node[ent, right=of mtr] (es)  {EvaluacionSaber};
\node[ent, right=of es]  (det) {DetalleEvalSaber};

\node[ent, below=of est] (ev) {Evaluacion};
\node[ent, right=of ev]  (de) {DetalleEval};
\node[ent, right=of de]  (at) {AlertaTematica};
\node[ent, right=of at]  (a)  {Alerta};

% relaciones
\draw[rel] (cu) -- (sec); \draw[rel] (al) -- (per); \draw[rel] (sec) -- (mtr);
\draw[rel] (est) -- (mtr); \draw[rel] (us) -- (sec); \draw[rel] (mat) -- (et);
\draw[rel] (ts) -- (et); \draw[rel] (et) -- (det); \draw[rel] (es) -- (det);
\draw[rel] (est) -- (es); \draw[rel] (us) -- (es); \draw[rel] (per) -- (es);
\draw[rel] (mat) -- (es); \draw[rel] (et) -- (at); \draw[rel] (mat) -- (at);
\draw[rel] (per) -- (at); \draw[rel] (est) -- (at); \draw[rel] (est) -- (ev);
\draw[rel] (per) -- (ev); \draw[rel] (us) -- (ev); \draw[rel] (ev) -- (de);
\draw[rel] (est) -- (a);
\end{tikzpicture}
\end{center}

\textbf{Entidades principales (27 totales):}
\code{Usuario}, \code{Rol}, \code{TokenRefresh},
\code{CentroEducativo}, \code{Nivel}, \code{Seccion}, \code{UsuarioSeccion},
\code{AnioLectivo}, \code{Periodo}, \code{Materia},
\code{Estudiante}, \code{Matricula},
\code{Evaluacion}, \code{DetalleEvaluacion}, \code{DimensionEvaluacion},
\code{CriterioIndicador}, \code{Contenido}, \code{EscalaValoracion},
\code{TipoSaber}, \code{EjeTematico}, \code{NivelDesempeno},
\code{EvaluacionSaber}, \code{DetalleEvaluacionSaber},
\code{Alerta}, \code{AlertaTematica}, \code{ConfiguracionAlerta},
\code{RegistroAuditoria}.

\subsection{Auditoría a nivel de base de datos}

Toda la auditoría está implementada \emph{exclusivamente en la BD}. La tabla
\code{registro\_auditoria} almacena documentos JSONB con los cambios; los \emph{triggers}
\code{fn\_actualizar\_updated\_at} mantienen las marcas temporales actualizadas. \textbf{No se
requiere ni se permite código de auditoría a nivel de aplicación.}

\subsection{Conjunto de datos semilla}

Los datos sembrados por V2–V11 garantizan reproducibilidad de pruebas y demostraciones:

\begin{table}[h]
\centering\renewcommand{\arraystretch}{1.15}\small
\begin{tabular}{@{}p{4.5cm} p{8.5cm}@{}}
\toprule
\textbf{Recurso} & \textbf{Identificadores semilla} \\
\midrule
Usuarios               & admin=1, mgarcia=2, kcortes=5 \\
Roles                  & ADMIN, COORDINADOR, DOCENTE \\
Secciones de prueba    & 1°A Bolívar = 1, 4°A Keylor = 4 \\
Periodo                & I Trimestre = 1 \\
Materias               & Matemáticas = 1, Español = 2, Ciencias = 3, Estudios Sociales = 4 \\
Ejes Matemáticas Conc. & ids 22--28 (7 ejes Conceptuales) \\
Ejes Español Conc.     & ids 1--7 \\
Estudiantes 4°A Keylor & ids 31--33 \\
Estudiantes 1°A María  & id 1 \\
\bottomrule
\end{tabular}
\caption{Identificadores semilla utilizados por las pruebas funcionales.}
\end{table}

% =====================================================================
\section{Módulo 1 — Autenticación y Seguridad (JWT)}
% =====================================================================

\subsection{Descripción funcional}

Este módulo gestiona la autenticación de usuarios mediante \textbf{JSON Web Tokens (JWT)}.
Implementa un esquema con \emph{access token} de corta duración y \emph{refresh token}
rotativo, y protege todos los endpoints excepto los públicos y el \emph{health check}.

\subsection{Endpoints}

\begin{table}[h]
\centering\renewcommand{\arraystretch}{1.2}\small
\begin{tabular}{@{}p{1.4cm} p{4.0cm} p{6.5cm} p{1.5cm}@{}}
\toprule
\textbf{Método} & \textbf{Ruta} & \textbf{Descripción} & \textbf{Auth} \\
\midrule
POST & \code{/api/auth/login}    & Autentica usuario; retorna access token + refresh token & No \\
POST & \code{/api/auth/refresh}  & Renueva el access token (rota el refresh)               & No \\
POST & \code{/api/auth/logout}   & Revoca el refresh token activo                          & No \\
GET  & \code{/api/auth/me}       & Retorna los datos del usuario autenticado               & Sí \\
\bottomrule
\end{tabular}
\caption{Endpoints expuestos por \code{AuthController}.}
\end{table}

\subsection{Contrato de la petición de \emph{login}}

\begin{lstlisting}[style=json,caption={POST \code{/api/auth/login} — \emph{Request body}}]
{
  "correo":   "admin@atara.mep.go.cr",   // @NotBlank @Email
  "password": "Admin1234!"               // @NotBlank
}
\end{lstlisting}

\begin{lstlisting}[style=json,caption={POST \code{/api/auth/login} — \emph{Response body} (200 OK)}]
{
  "accessToken":  "<JWT>",
  "refreshToken": "<UUID>",
  "tokenType":    "Bearer",
  "expiresIn":    3600,
  "userId":       1,
  "correo":       "admin@atara.mep.go.cr",
  "nombre":       "Administrador",
  "apellidos":    "Sistema",
  "rol":          "ADMIN"
}
\end{lstlisting}

\subsection{Reglas de seguridad}

\begin{ataraInfo}
\begin{itemize}[leftmargin=*]
  \item Las contraseñas se almacenan con \textbf{BCryptPasswordEncoder} (\emph{strength} 12).
        Nunca se exponen en texto plano ni se registran en logs.
  \item El \emph{access token} se valida en \code{JwtAuthenticationFilter} en cada request,
        nunca dentro de controladores ni servicios.
  \item El \emph{refresh token} \textbf{rota} en cada uso: el anterior se revoca y se emite
        uno nuevo.
  \item Todos los endpoints exigen el header \code{Authorization: Bearer <token>}, excepto
        \code{/api/auth/login}, \code{/api/auth/refresh}, \code{/api/auth/logout} y
        \code{/actuator/health}.
  \item La sesión es \textbf{stateless}: \code{SessionCreationPolicy.STATELESS}.
  \item \emph{CSRF deshabilitado}, \emph{CORS configurable} mediante
        \code{CORS\_ALLOWED\_ORIGINS}.
\end{itemize}
\end{ataraInfo}

\subsection{Roles disponibles}

\begin{table}[h]
\centering\renewcommand{\arraystretch}{1.2}\small
\begin{tabular}{@{}p{2.8cm} p{4.5cm} p{6.2cm}@{}}
\toprule
\textbf{Rol} & \textbf{Descripción} & \textbf{Acceso} \\
\midrule
\code{ADMIN}        & Administrador del sistema   & Total (todas las secciones y materias) \\
\code{COORDINADOR}  & Coordinador académico       & Total (equivalente a ADMIN) \\
\code{DOCENTE}      & Docente de aula             & Filtrado por secciones y materias asignadas \\
\bottomrule
\end{tabular}
\caption{Roles del sistema y alcance de su acceso.}
\end{table}

\subsection{Manejo de respuestas no autorizadas}

\code{SecurityConfig} registra dos manejadores de error con cuerpo JSON consistente con el
\code{ApiResponse} estándar:

\begin{itemize}[leftmargin=*]
  \item \textbf{401 Unauthorized} — \code{authenticationEntryPoint}: \emph{“Token de acceso
        requerido o inválido”}.
  \item \textbf{403 Forbidden} — \code{accessDeniedHandler}: \emph{“No tiene permisos para
        acceder a este recurso”}.
\end{itemize}

\subsection{Evidencia de pruebas funcionales}

\begin{table}[h]
\centering\renewcommand{\arraystretch}{1.15}\small
\begin{tabular}{@{}p{7.0cm} p{5.0cm} p{1.5cm}@{}}
\toprule
\textbf{Caso de prueba} & \textbf{Resultado esperado} & \textbf{Estado} \\
\midrule
Login con credenciales correctas (admin)    & 200 OK + accessToken presente    & \pass \\
Login de docente                            & 200 OK + rol = DOCENTE           & \pass \\
Contraseña incorrecta                       & 401 Unauthorized                 & \pass \\
Correo inexistente                          & 401 Unauthorized                 & \pass \\
Formato de correo inválido                  & 400 Bad Request                  & \pass \\
Cuerpo vacío en login                       & 400 Bad Request                  & \pass \\
Endpoint protegido sin token                & 401 Unauthorized                 & \pass \\
Endpoint protegido con token inválido       & 401 Unauthorized                 & \pass \\
GET \code{/api/auth/me} con token válido    & 200 OK + datos del usuario       & \pass \\
GET \code{/api/auth/me} sin token           & 401 Unauthorized                 & \pass \\
\midrule
\textbf{Total}                              &                                  & \textbf{10 / 10} \\
\bottomrule
\end{tabular}
\end{table}

% =====================================================================
\section{Módulo 2 — Control de Acceso por Rol (\code{ContextoUsuario})}
% =====================================================================

\subsection{Descripción funcional}

El control de acceso basado en rol se implementa mediante el patrón \textbf{ContextoUsuario}:
un \code{record} inmutable que encapsula el alcance del usuario autenticado. Se construye
\emph{una sola vez por request} en \code{ContextoUsuarioService} y se pasa como parámetro a
los servicios que necesitan verificar permisos.

\begin{ataraInfo}
\textbf{Punto único de acceso:} \code{ContextoUsuarioService.obtenerContextoActual()} es la
\textbf{única} clase que lee \code{SecurityContextHolder}. Los servicios de dominio nunca
acceden directamente al contexto de seguridad de Spring.
\end{ataraInfo}

\subsection{Estructura del \code{ContextoUsuario}}

\begin{lstlisting}[style=java,caption={\code{ContextoUsuario.java} — \emph{record} inmutable por request}]
public record ContextoUsuario(
    Long          usuarioId,
    RolNombre     rol,
    boolean       esAdmin,        // true para ADMIN y COORDINADOR
    Set<Long>     seccionIds,     // IDs de secciones asignadas (solo docentes)
    Set<Integer>  materiaIds      // IDs de materias asignadas (solo docentes)
) {
    // Lanza AccesoDenegadoException (-> HTTP 403) si no tiene acceso
    public void verificarSeccion(Long seccionId)  { ... }
    public void verificarMateria(Integer materiaId) { ... }
    public boolean tieneSeccion(Long seccionId)   { ... }
    public boolean tieneMateria(Integer materiaId){ ... }
}
\end{lstlisting}

\subsection{Flujo de construcción del contexto}

\begin{enumerate}[leftmargin=*]
  \item \code{JwtAuthenticationFilter} valida el token y carga el \code{UserDetails} en el
        \code{SecurityContextHolder}.
  \item El controller llama a \code{ContextoUsuarioService.obtenerContextoActual()}.
  \item El servicio lee \code{SecurityContextHolder}, consulta las tablas
        \code{usuario\_secciones} y \code{usuario\_materias}, e instancia el
        \code{ContextoUsuario}.
  \item Los servicios de dominio reciben el contexto como parámetro y llaman a
        \code{contexto.verificarSeccion(id)} / \code{contexto.verificarMateria(id)} antes de
        operar.
  \item Si la verificación falla, se lanza \code{AccesoDenegadoException} (mapeado a 403
        por \code{GlobalExceptionHandler}).
\end{enumerate}

\subsection{Filtrado automático de catálogos por rol}

El catálogo de materias (\code{GET /api/catalogos/saberes/materias}) aplica el filtro de
manera automática:

\begin{itemize}[leftmargin=*]
  \item \textbf{ADMIN / COORDINADOR}: devuelve las cinco materias del sistema.
  \item \textbf{DOCENTE}: devuelve únicamente las materias contenidas en su \code{materiaIds}.
\end{itemize}

\subsection{Evidencia de pruebas funcionales}

\begin{table}[h]
\centering\renewcommand{\arraystretch}{1.15}\small
\begin{tabular}{@{}p{6.5cm} p{5.5cm} p{1.5cm}@{}}
\toprule
\textbf{Caso de prueba} & \textbf{Resultado esperado} & \textbf{Estado} \\
\midrule
Admin ve catálogo de materias               & 5 materias                       & \pass \\
Docente Keylor ve materias                  & Solo Matemáticas (1 materia)     & \pass \\
Coordinador ve materias                     & 5 materias (acceso total)        & \pass \\
Catálogo de materias sin token              & 401 Unauthorized                 & \pass \\
Catálogo de tipos de saber                  & 3 tipos (sin filtro por rol)     & \pass \\
Catálogo de niveles de desempeño            & 5 niveles                        & \pass \\
Ejes de Matemáticas Conceptual              & 7 ejes                           & \pass \\
Ejes de Español Conceptual                  & 7 ejes                           & \pass \\
Secciones del docente Keylor                & $\geq$ 1 sección (4°A)            & \pass \\
Secciones de la docente María               & $\geq$ 2 secciones                & \pass \\
\midrule
\textbf{Total}                              &                                  & \textbf{10 / 10} \\
\bottomrule
\end{tabular}
\end{table}

% =====================================================================
\section{Módulo 3 — Gestión de Estudiantes y Matrículas}
% =====================================================================

\subsection{Descripción funcional}

Este módulo administra el registro y modificación de estudiantes, así como la inscripción
(matrícula) de cada estudiante en una sección y un año lectivo. Se diseñó pensando en la
reutilización: la creación es \emph{idempotente} por número de cédula, lo que permite
re-importar listas (por ejemplo, desde el módulo PIAD) sin generar conflictos de unicidad.

\subsection{Endpoints}

\begin{table}[h]
\centering\renewcommand{\arraystretch}{1.15}\small
\begin{tabular}{@{}p{1.4cm} p{6.0cm} p{6.5cm}@{}}
\toprule
\textbf{Método} & \textbf{Ruta} & \textbf{Descripción} \\
\midrule
POST & \code{/api/estudiantes}                                  & Crea o recupera un estudiante (idempotente por cédula) \\
GET  & \code{/api/estudiantes}                                  & Lista todos los estudiantes \\
GET  & \code{/api/estudiantes/\{id\}}                           & Consulta por ID \\
PUT  & \code{/api/estudiantes/\{id\}}                           & Actualiza datos del estudiante \\
\midrule
POST & \code{/api/matriculas}                                   & Registra matrícula \\
GET  & \code{/api/matriculas/estudiante/\{estudianteId\}}       & Lista matrículas de un estudiante \\
GET  & \code{/api/matriculas/seccion/\{seccionId\}}             & Lista matrículas de una sección \\
\bottomrule
\end{tabular}
\caption{Endpoints de \code{EstudianteController} y \code{MatriculaController}.}
\end{table}

\subsection{Reglas de negocio}

\begin{ataraInfo}
\begin{itemize}[leftmargin=*]
  \item \textbf{Idempotencia por cédula:} \code{POST /api/estudiantes} ya \emph{no} lanza
        409 Conflict cuando la cédula ya existe; devuelve el estudiante existente.
  \item \textbf{Validación:} la cédula debe cumplir el formato establecido por las
        anotaciones \code{@NotBlank} y \code{@Pattern} en
        \code{EstudianteRequestDto}.
  \item \textbf{Restricciones de matrícula:} no se permite matrícula doble del mismo
        estudiante en el mismo año lectivo y sección activa.
\end{itemize}
\end{ataraInfo}

% =====================================================================
\section{Módulo 4 — Años Lectivos y Periodos}
% =====================================================================

\subsection{Descripción funcional}

Gestiona la estructura temporal del sistema: años lectivos y los periodos (trimestres)
que contienen. Al crear un año lectivo, el servicio genera automáticamente tres trimestres
dividiendo el rango de fechas en partes iguales y dejando el primero como activo.

\subsection{Endpoints}

\begin{table}[h]
\centering\renewcommand{\arraystretch}{1.15}\small
\begin{tabular}{@{}p{1.4cm} p{6.5cm} p{6.0cm}@{}}
\toprule
\textbf{Método} & \textbf{Ruta} & \textbf{Descripción} \\
\midrule
POST   & \code{/api/anios-lectivos}                       & Crea un año lectivo y genera 3 trimestres \\
GET    & \code{/api/anios-lectivos}                       & Lista todos los años \\
GET    & \code{/api/anios-lectivos/activo}                & Año actualmente activo \\
GET    & \code{/api/anios-lectivos/\{id\}}                & Consulta por ID \\
PUT    & \code{/api/anios-lectivos/\{id\}}                & Actualiza un año \\
PUT    & \code{/api/anios-lectivos/\{id\}/activar}        & Activa este año (desactiva el resto) \\
DELETE & \code{/api/anios-lectivos/\{id\}}                & Elimina un año \\
\midrule
POST   & \code{/api/periodos}                             & Crea un periodo manualmente \\
GET    & \code{/api/periodos?anioLectivoId=}              & Lista periodos del año \\
GET    & \code{/api/periodos/activo?anioLectivoId=}       & Periodo activo del año \\
PUT    & \code{/api/periodos/\{id\}}                      & Actualiza un periodo \\
PUT    & \code{/api/periodos/\{id\}/activar}              & Activa el periodo (desactiva el resto del año) \\
DELETE & \code{/api/periodos/\{id\}}                      & Elimina un periodo \\
\bottomrule
\end{tabular}
\caption{Endpoints de \code{AnioLectivoController} y \code{PeriodoController}.}
\end{table}

\subsection{Reglas de negocio}

\begin{ataraInfo}
\begin{itemize}[leftmargin=*]
  \item \textbf{Auto-creación de periodos:} \code{AnioLectivoServiceImpl.crearAnioLectivo()}
        divide el rango \code{[fechaInicio, fechaFin]} en tres partes iguales y crea los
        trimestres I, II y III. El I queda activo.
  \item \textbf{Activación exclusiva:} activar un año o un periodo desactiva todos los demás
        del mismo ámbito; el sistema garantiza que solo exista uno activo a la vez.
  \item \textbf{Integridad referencial:} no se permite eliminar un año lectivo si tiene
        matrículas o evaluaciones asociadas.
\end{itemize}
\end{ataraInfo}

% =====================================================================
\section{Módulo 5 — Secciones y Docentes}
% =====================================================================

\subsection{Descripción funcional}

El módulo de secciones gestiona la estructura académica del sistema: agrupaciones de
estudiantes asignadas a un docente para un año lectivo. Expone catálogos auxiliares
(niveles, centros, docentes) que permiten al frontend construir formularios sin depender
de datos codificados.

\subsection{Endpoints}

\begin{table}[h]
\centering\renewcommand{\arraystretch}{1.15}\small
\begin{tabular}{@{}p{1.6cm} p{6.5cm} p{5.5cm}@{}}
\toprule
\textbf{Método} & \textbf{Ruta} & \textbf{Descripción} \\
\midrule
GET    & \code{/api/secciones?anioLectivoId=}             & Lista secciones del año (filtrada por rol) \\
GET    & \code{/api/secciones/docente/\{docenteId\}}      & Secciones asignadas a un docente \\
POST   & \code{/api/secciones}                            & Crea una nueva sección \\
PUT    & \code{/api/secciones/\{id\}}                     & Actualiza una sección \\
DELETE & \code{/api/secciones/\{id\}}                     & Elimina una sección \\
GET    & \code{/api/secciones/catalogos/niveles}          & Catálogo de niveles educativos \\
GET    & \code{/api/secciones/catalogos/centros}          & Catálogo de centros educativos \\
GET    & \code{/api/secciones/catalogos/docentes}         & Catálogo de docentes disponibles \\
\bottomrule
\end{tabular}
\caption{Endpoints de \code{SeccionController}.}
\end{table}

\subsection{Datos de control de acceso}

La tabla \code{usuario\_secciones} relaciona docentes con sus secciones asignadas; la tabla
\code{usuario\_materias} relaciona docentes con sus materias asignadas. Ambas son consultadas
por \code{ContextoUsuarioService} al construir el \code{ContextoUsuario} de cada request.

\begin{table}[h]
\centering\renewcommand{\arraystretch}{1.15}\small
\begin{tabular}{@{}p{4.7cm} p{2.8cm} p{6.0cm}@{}}
\toprule
\textbf{Usuario} & \textbf{Rol} & \textbf{Secciones / materias asignadas} \\
\midrule
\code{admin@atara.mep.go.cr}     & ADMIN        & Todas (sin restricción) \\
\code{jperez@atara.mep.go.cr}    & COORDINADOR  & Todas (acceso total) \\
\code{mgarcia@atara.mep.go.cr}   & DOCENTE      & Secciones 1 y 3; todas las materias \\
\code{kcortes@atara.mep.go.cr}   & DOCENTE      & Sección 4; solo Matemáticas (id = 1) \\
\code{avargas@atara.mep.go.cr}   & DOCENTE      & Materias adicionales: Español y Ciencias (V10) \\
\bottomrule
\end{tabular}
\caption{Usuarios semilla y su alcance configurado.}
\end{table}

\begin{ataraInfo}
Todas las contraseñas semilla son \code{Admin1234!}. En producción, los \emph{hashes} BCrypt
de \code{V2\_\_sample\_data.sql} deben ser reemplazados por contraseñas reales generadas con
\code{new BCryptPasswordEncoder(12).encode(...)}.
\end{ataraInfo}

% =====================================================================
\section{Módulo 6 — Catálogos de Saberes}
% =====================================================================

\subsection{Descripción funcional}

Los catálogos de saberes exponen la información de referencia que el frontend necesita para
construir los formularios de evaluación: tipos de saber, materias filtradas por rol, ejes
temáticos por materia y tipo, y niveles de desempeño.

\subsection{Endpoints}

\begin{table}[h]
\centering\renewcommand{\arraystretch}{1.2}\small
\begin{tabular}{@{}p{1.4cm} p{6.5cm} p{6.0cm}@{}}
\toprule
\textbf{Método} & \textbf{Ruta} & \textbf{Descripción} \\
\midrule
GET & \code{/api/catalogos/saberes/tipos}                       & Lista los 3 tipos de saber \\
GET & \code{/api/catalogos/saberes/materias}                    & Lista materias (filtradas por rol del usuario) \\
GET & \code{/api/catalogos/saberes/ejes?materiaId=\&tipoSaberId=} & Lista ejes (parámetros opcionales) \\
GET & \code{/api/catalogos/saberes/niveles-desempeno}           & Lista los 5 niveles de desempeño \\
\bottomrule
\end{tabular}
\caption{Endpoints de \code{CatalogoSaberController}.}
\end{table}

\subsection{Datos de referencia}

\textbf{Tipos de saber (3)}

\begin{table}[h]\centering\small
\begin{tabular}{@{}c l l@{}}
\toprule
\textbf{ID} & \textbf{Clave} & \textbf{Nombre} \\
\midrule
1 & \code{CONCEPTUAL}     & Saber Conceptual    \\
2 & \code{PROCEDIMENTAL}  & Saber Procedimental \\
3 & \code{ACTITUDINAL}    & Saber Actitudinal   \\
\bottomrule
\end{tabular}
\end{table}

\textbf{Niveles de desempeño (5)}

\begin{table}[h]\centering\small
\begin{tabular}{@{}c l l@{}}
\toprule
\textbf{Valor} & \textbf{Nombre} & \textbf{Descripción breve} \\
\midrule
1 & Inicial       & No evidencia dominio \\
2 & En desarrollo & Dominio parcial \\
3 & Intermedio    & Dominio básico \\
4 & Logrado       & Dominio sólido \\
5 & Avanzado      & Dominio pleno con transferencia \\
\bottomrule
\end{tabular}
\end{table}

\textbf{Materias y conteo de ejes}

\begin{table}[h]\centering\small
\begin{tabular}{@{}c l l@{}}
\toprule
\textbf{ID} & \textbf{Nombre} & \textbf{Ejes totales} \\
\midrule
1 & Español           & 21 (7 por tipo de saber) \\
2 & Matemáticas       & 21 (7 por tipo de saber) \\
3 & Ciencias          & 21 (7 por tipo de saber) \\
4 & Estudios Sociales & 21 (7 por tipo de saber) \\
5 & (adicional)       & --- \\
\bottomrule
\end{tabular}
\end{table}

% =====================================================================
\section{Módulo 7 — Evaluaciones por Saber}
% =====================================================================

\subsection{Descripción funcional}

El módulo de \textbf{Evaluaciones por Saber} permite al docente registrar la valoración de
un estudiante en cada eje temático de una materia, para un periodo lectivo específico.
Cada evaluación cubre un \emph{tipo de saber} (Conceptual, Procedimental o Actitudinal) y
contiene un detalle por cada eje seleccionado, en escala 1--5.

\subsection{Endpoints}

\begin{table}[h]
\centering\renewcommand{\arraystretch}{1.15}\small
\begin{tabularx}{\textwidth}{@{}p{1.4cm} X p{3.8cm}@{}}
\toprule
\textbf{Método} & \textbf{Ruta} & \textbf{Descripción} \\
\midrule
POST & \ruta{/api/evaluaciones-saber}                                              & Registra nueva evaluación \\
PUT  & \ruta{/api/evaluaciones-saber/{id}}                                         & Recalifica (reemplaza detalles) \\
GET  & \ruta{/api/evaluaciones-saber/{id}}                                         & Consulta por ID \\
GET  & \ruta{/api/evaluaciones-saber/estudiante/{eId}/periodo/{pId}}               & Lista por estudiante y periodo \\
GET  & \ruta{/api/evaluaciones-saber/seccion/{sId}/periodo/{pId}}                  & Lista por sección y periodo \\
GET  & \ruta{/api/evaluaciones-saber/promedios/estudiante/{eId}/periodo/{pId}}     & Promedios del estudiante \\
GET  & \ruta{/api/evaluaciones-saber/promedios/seccion/{sId}/periodo/{pId}}        & Promedios de la sección \\
\bottomrule
\end{tabularx}
\caption{Endpoints de \code{EvaluacionSaberController}.}
\end{table}

\subsection{Contrato de la petición}

\begin{lstlisting}[style=json,caption={POST \code{/api/evaluaciones-saber} — \emph{Request body}}]
{
  "estudianteId":  31,
  "periodoId":     1,
  "usuarioId":     5,
  "seccionId":     4,
  "materiaId":     1,
  "tipoSaberId":   1,
  "fechaEvaluacion": "2026-04-15",   // opcional
  "observacion":     "Observacion libre",   // opcional
  "detalles": [
    { "ejeTemaaticoId": 22, "valor": 3 },
    { "ejeTemaaticoId": 23, "valor": 4 },
    { "ejeTemaaticoId": 24, "valor": 3 }
    // ... un objeto por eje evaluado
  ]
}
\end{lstlisting}

\begin{ataraNota}
\textbf{Erratographic intencional:} el campo se llama \code{ejeTemaaticoId} (con doble «a»)
en todo el sistema —DTO, generador JSON de \code{IntegrationTestBase} y frontend— y se
mantiene de manera deliberada para no introducir cambios coordinados frágiles. \emph{Cualquier
mantenimiento debe respetar la grafía existente.}
\end{ataraNota}

\begin{lstlisting}[style=json,caption={POST \code{/api/evaluaciones-saber} — \emph{Response body} (201 Created)}]
{
  "id":                      101,
  "estudianteId":            31,
  "estudianteNombreCompleto": "Andrea Mora Vega",
  "periodoId":               1,
  "periodoNombre":           "I Trimestre 2026",
  "materiaId":               1,
  "materiaNombre":           "Matematicas",
  "tipoSaberId":             1,
  "tipoSaberNombre":         "Saber Conceptual",
  "seccionId":               4,
  "seccionNombre":           "4 grado A",
  "detalles": [
    { "ejeTemaaticoId": 22, "ejeNombre": "...", "valor": 3 }
  ],
  "createdAt": "2026-04-15T09:00:00"
}
\end{lstlisting}

\subsection{Escala de valoración por saber}

\begin{table}[h]
\centering\renewcommand{\arraystretch}{1.15}\small
\begin{tabular}{@{}c l p{8.5cm}@{}}
\toprule
\textbf{Valor} & \textbf{Nivel} & \textbf{Descripción} \\
\midrule
1 & Inicial       & No evidencia dominio del eje \\
2 & En desarrollo & Dominio parcial, requiere apoyo constante \\
3 & Intermedio    & Dominio básico con algunos errores \\
4 & Logrado       & Dominio sólido con mínimos errores \\
5 & Avanzado      & Dominio pleno y transferencia de conocimiento \\
\bottomrule
\end{tabular}
\caption{Escala de 5 puntos aplicada a las evaluaciones por saber.}
\end{table}

\subsection{Reglas de negocio}

\begin{ataraInfo}
\begin{itemize}[leftmargin=*]
  \item \textbf{Restricción de unicidad:} una sola evaluación por
        \code{(estudiante\_id, usuario\_id, periodo\_id)} (definida a nivel de BD).
  \item \textbf{Validaciones de entrada:} \code{detalles} debe tener al menos un elemento
        (\code{@Size(min=1)}); \code{valor} debe estar entre 1 y 5
        (\code{@Min(1)}, \code{@Max(5)}).
  \item \textbf{Recalificación (PUT):} reemplaza todos los detalles existentes mediante
        \code{orphanRemoval=true} y luego inserta los nuevos. \emph{No} crea una segunda
        evaluación.
  \item \textbf{Control de acceso:} el servicio verifica que el docente tenga acceso a la
        sección y a la materia antes de registrar.
\end{itemize}
\end{ataraInfo}

\subsection{Ejes temáticos por materia}

Cada materia posee \textbf{21 ejes temáticos independientes} (7 por tipo de saber):

\begin{table}[h]\centering\small
\begin{tabular}{@{}l l c l@{}}
\toprule
\textbf{Materia} & \textbf{Tipo de saber} & \textbf{Ejes} & \textbf{IDs (ejemplo)} \\
\midrule
\multirow{3}{*}{Español}     & Conceptual    & 7 & 1--7   \\
                             & Procedimental & 7 & 8--14  \\
                             & Actitudinal   & 7 & 15--21 \\ \midrule
\multirow{3}{*}{Matemáticas} & Conceptual    & 7 & 22--28 \\
                             & Procedimental & 7 & 29--35 \\
                             & Actitudinal   & 7 & 36--42 \\ \midrule
Ciencias, Est. Sociales      & Cada uno 3 tipos & 7 c/u & (V6) \\
\bottomrule
\end{tabular}
\caption{Distribución de ejes temáticos por materia.}
\end{table}

\subsection{Evidencia de pruebas funcionales}

\begin{table}[h]\centering\small\renewcommand{\arraystretch}{1.15}
\begin{tabular}{@{}p{6.5cm} p{5.5cm} p{1.4cm}@{}}
\toprule
\textbf{Caso de prueba} & \textbf{Resultado esperado} & \textbf{Estado} \\
\midrule
Registrar evaluación válida                       & 201 Created + ID generado                & \pass \\
Registrar y recuperar por ID                      & Mismo ID y detalles coherentes           & \pass \\
Registrar y consultar promedios                   & 200 OK + \code{promediosPorTipoSaber} presente & \pass \\
Listar por estudiante y periodo                   & Lista con $\geq$ 1 evaluación             & \pass \\
Registrar sin token                               & 401 Unauthorized                         & \pass \\
Registrar con detalles vacíos                     & 400 Bad Request                          & \pass \\
Registrar con valor fuera de rango (6)            & 400 Bad Request                          & \pass \\
Buscar por ID inexistente (999999)                & 404 Not Found                            & \pass \\
\midrule
\textbf{Total}                                    &                                          & \textbf{8 / 8} \\
\bottomrule
\end{tabular}
\end{table}

% =====================================================================
\section{Módulo 8 — Alertas Temáticas}
% =====================================================================

\subsection{Descripción funcional}

El módulo de \textbf{Alertas Temáticas} es el núcleo del sistema de alerta temprana. Analiza
las evaluaciones registradas para un estudiante (o una sección completa) en un periodo,
calcula el promedio por eje temático y genera automáticamente alertas para los ejes cuyo
promedio no alcanza el umbral de desempeño esperado.

\subsection{Endpoints}

\begin{table}[h]
\centering\renewcommand{\arraystretch}{1.15}\small
\begin{tabularx}{\textwidth}{@{}p{1.4cm} X p{4.0cm}@{}}
\toprule
\textbf{Método} & \textbf{Ruta} & \textbf{Descripción} \\
\midrule
POST & \ruta{/api/alertas-tematicas/generar/estudiante/{eId}/periodo/{pId}} & Genera alertas para un estudiante \\
POST & \ruta{/api/alertas-tematicas/generar/seccion/{sId}/periodo/{pId}}    & Genera alertas para toda la sección \\
GET  & \ruta{/api/alertas-tematicas/estudiante/{eId}/periodo/{pId}}         & Consulta alertas de un estudiante \\
GET  & \ruta{/api/alertas-tematicas/seccion/{sId}/periodo/{pId}}            & Consulta alertas de una sección \\
\bottomrule
\end{tabularx}
\caption{Endpoints de \code{AlertaTematicaController}.}
\end{table}

\subsection{Respuesta de alerta}

\begin{lstlisting}[style=json,caption={Ejemplo de alerta generada (nivel ALTA)}]
{
  "id":                      15,
  "estudianteId":            32,
  "estudianteNombreCompleto":"Luis Perez Mora",
  "periodoId":               1,
  "periodoNombre":           "I Trimestre 2026",
  "ejeTemaaticoId":          22,
  "ejeNombre":               "Algebra y pensamiento algebraico",
  "ejeClave":                "MAT_C_ALGEBRA",
  "materiaId":               1,
  "materiaNombre":           "Matematicas",
  "tipoSaberId":             1,
  "tipoSaberNombre":         "Saber Conceptual",
  "promedio":                1.43,
  "nivelAlerta":             "ALTA",
  "motivo":                  "Luis Perez Mora presenta promedio 1.43 en Matematicas ...",
  "estado":                  "ACTIVA",
  "fechaGeneracion":         "2026-04-15T09:05:00"
}
\end{lstlisting}

\subsection{Reglas de generación de alertas}

\begin{table}[h]
\centering\renewcommand{\arraystretch}{1.2}\small
\begin{tabular}{@{}c l p{8.3cm}@{}}
\toprule
\textbf{Promedio} & \textbf{Nivel} & \textbf{Interpretación} \\
\midrule
$\leq 2{,}0$         & {\color{ataraRojo}\textbf{ALTA}}        & Riesgo crítico. Requiere intervención inmediata. \\
$2{,}1 - 3{,}0$      & {\color{ataraNaranja}\textbf{MEDIA}}    & Riesgo moderado. Requiere seguimiento. \\
$> 3{,}0$            & (SIN\_ALERTA)                           & Desempeño aceptable. No se genera ni persiste alerta. \\
\bottomrule
\end{tabular}
\caption{Niveles de alerta según el promedio de eje temático.}
\end{table}

\textbf{Reglas adicionales:}

\begin{itemize}[leftmargin=*]
  \item Las alertas con nivel \code{SIN\_ALERTA} \textbf{no se persisten} en la base de datos.
  \item Si ya existe una alerta para la tupla \code{(estudiante, periodo, eje, materia)}, la
        nueva generación la \textbf{reemplaza}.
  \item El estado inicial de toda alerta generada es \code{ACTIVA}.
  \item La generación por sección delega internamente en la generación por estudiante para
        cada estudiante matriculado.
\end{itemize}

\subsection{Evidencia de pruebas funcionales}

\begin{table}[h]\centering\small\renewcommand{\arraystretch}{1.15}
\begin{tabular}{@{}p{6.5cm} p{5.5cm} p{1.4cm}@{}}
\toprule
\textbf{Caso de prueba} & \textbf{Resultado esperado} & \textbf{Estado} \\
\midrule
Valores 1 (promedio $\leq$ 2) generan alerta       & Nivel = ALTA                              & \pass \\
Valores 5 (promedio $>$ 3) no generan alertas      & Lista vacía (\code{SIN\_ALERTA} no se persiste) & \pass \\
Generar y consultar — cantidad coherente           & GET retorna misma cantidad que POST       & \pass \\
Generar alertas sin token                          & 401 Unauthorized                          & \pass \\
Estudiante sin evaluaciones                        & Lista vacía (200 OK)                      & \pass \\
\midrule
\textbf{Total}                                     &                                           & \textbf{5 / 5} \\
\bottomrule
\end{tabular}
\end{table}

% =====================================================================
\section{Módulo 9 — Evaluaciones (sistema 4-puntos) y Alertas Tempranas}
% =====================================================================

\subsection{Descripción funcional}

El sistema mantiene además el modelo original de \emph{evaluaciones por dimensiones} (escala
4-puntos) que precedió al modelo de saberes. Persiste para soportar los reportes históricos y
las alertas tempranas (\code{AlertaController}).

\subsection{Escala de valoración (original)}

Escala de 4 puntos sobre 5 dimensiones de evaluación: Insuficiente (1), Básico (2),
Satisfactorio (3), Destacado (4). Las dimensiones incluyen Rendimiento Académico,
Participación, Hábitos de Estudio y Factores Socioemocionales.

\subsection{Endpoints}

\begin{table}[h]
\centering\renewcommand{\arraystretch}{1.15}\small
\begin{tabularx}{\textwidth}{@{}p{1.4cm} X p{4.2cm}@{}}
\toprule
\textbf{Método} & \textbf{Ruta} & \textbf{Descripción} \\
\midrule
POST & \ruta{/api/evaluaciones}                                                    & Crea evaluación 4-puntos \\
POST & \ruta{/api/evaluaciones/{id}/detalles}                                      & Agrega detalles a evaluación \\
GET  & \ruta{/api/evaluaciones/{id}}                                               & Consulta por ID \\
GET  & \ruta{/api/evaluaciones/estudiante/{estudianteId}}                          & Lista por estudiante \\
GET  & \ruta{/api/evaluaciones/estudiante/{estudianteId}/periodo/{periodoId}}      & Por estudiante y periodo \\
GET  & \ruta{/api/evaluaciones/periodo/{periodoId}}                                & Lista del periodo \\
\midrule
GET  & \ruta{/api/alertas/estudiante/{studentId}?periodoId=}                       & Alertas tempranas de un estudiante \\
GET  & \ruta{/api/alertas/seccion/{sectionId}?periodoId=}                          & Alertas tempranas de una sección \\
\bottomrule
\end{tabularx}
\caption{Endpoints de \code{EvaluacionController} y \code{AlertaController}.}
\end{table}

\textbf{Restricción única:} \code{(estudiante\_id, usuario\_id, periodo\_id)} sobre
\code{evaluaciones} (igual que en \code{evaluaciones\_saber}).

% =====================================================================
\section{Módulo 10 — Visualizaciones y Reportes}
% =====================================================================

\subsection{Visualizaciones}

\code{VisualizacionController} expone agregaciones diseñadas para alimentar las gráficas
de Chart.js en el frontend.

\begin{table}[h]
\centering\renewcommand{\arraystretch}{1.15}\small
\begin{tabularx}{\textwidth}{@{}p{1.4cm} X p{4.0cm}@{}}
\toprule
\textbf{Método} & \textbf{Ruta} & \textbf{Descripción} \\
\midrule
GET & \ruta{/api/visualizaciones/seccion/{sectionId}/distribucion}
    \newline\code{?materiaId=\&periodoId=}
    & Distribución de niveles de desempeño en una sección, materia y periodo. \\
\bottomrule
\end{tabularx}
\caption{Endpoints de \code{VisualizacionController}.}
\end{table}

DTOs principales: \code{DistribucionAlertasDto}, \code{GraficoValorDto}, \code{ProgresoSeccionDto}.

\subsection{Reportes}

\code{ReporteController} produce reportes consolidados por estudiante.

\begin{table}[h]
\centering\renewcommand{\arraystretch}{1.15}\small
\begin{tabular}{@{}p{1.4cm} p{7.5cm} p{6.0cm}@{}}
\toprule
\textbf{Método} & \textbf{Ruta} & \textbf{Descripción} \\
\midrule
GET & \code{/api/reportes/estudiante/\{studentId\}?materiaId=\&periodoId=}
    & Reporte completo del estudiante con detalles, promedios y alertas. \\
\bottomrule
\end{tabular}
\end{table}

DTOs principales: \code{ReporteEstudianteDto}, \code{ReporteSeccionDto},
\code{DetalleContenidoReporteDto}.

% =====================================================================
\section{Módulo 11 — Extracción PIAD (OCR)}
% =====================================================================

\subsection{Descripción funcional}

El módulo \textbf{PIAD} permite al docente cargar el PDF oficial de la Lista PIAD (\emph{Plan
Individual de Adecuación}), que el sistema procesa mediante OCR para extraer automáticamente
la nómina de estudiantes con adecuaciones. Esto elimina la necesidad de registrar manualmente
a estos estudiantes y agiliza la importación al inicio del año lectivo.

\subsection{Endpoint}

\begin{table}[h]
\centering\small\renewcommand{\arraystretch}{1.15}
\begin{tabular}{@{}p{1.4cm} p{4.5cm} p{8.5cm}@{}}
\toprule
\textbf{Método} & \textbf{Ruta} & \textbf{Descripción} \\
\midrule
POST & \code{/api/piad/extraer} & Recibe PDF \emph{multipart}, aplica OCR y retorna la lista
                                  de estudiantes extraída en JSON. \\
\bottomrule
\end{tabular}
\end{table}

\subsection{Tecnología de procesamiento}

\begin{enumerate}[leftmargin=*]
  \item \textbf{Apache PDFBox 3.0.3:} renderiza cada página del PDF a imagen \emph{raster}
        a 500 DPI (mejor reconocimiento de fuentes pequeñas).
  \item \textbf{Tess4j 5.11.0} (\emph{wrapper} de Tesseract): aplica OCR en español
        (\code{spa}) sobre las imágenes con PSM 7 (línea única) y \code{OcrEngineMode = 1}
        (LSTM, red neuronal).
  \item \textbf{Detección de bandas horizontales:} preprocesamiento que segmenta filas de la
        tabla por brillo medio, mejorando la precisión del OCR sobre tablas densas.
  \item \textbf{Parser por expresiones regulares:} extrae nombres, cédulas y tipos de
        adecuación a partir del texto crudo.
\end{enumerate}

\subsection{Estado de implementación}

\begin{ataraNota}
La extracción OCR está \textbf{implementada y operativa}. La \emph{persistencia} de los
estudiantes extraídos en la base de datos \textbf{no está implementada} en esta versión
(trabajo pendiente para V12). El endpoint retorna la lista como respuesta JSON pero no la
guarda. La idempotencia del endpoint \code{POST /api/estudiantes} permite que el frontend
re-importe la lista sin generar conflictos.
\end{ataraNota}

\subsection{Configuración requerida}

\begin{itemize}[leftmargin=*]
  \item Tesseract debe estar instalado con el paquete de idioma español (\code{spa}).
  \item Configurar \code{piad.tessdata.dir} en \code{application.properties} o establecer la
        variable de entorno \code{TESSDATA\_PREFIX}.
  \item Si ninguna ruta es válida, el servicio prueba ubicaciones estándar de Linux:
        \code{/usr/share/tesseract-ocr/5/tessdata}, \code{/usr/share/tesseract-ocr/4/tessdata},
        \code{/usr/local/share/tessdata}.
\end{itemize}

\textbf{Limitaciones conocidas} (ver \code{docs/ocr-pdfbox-tess4j.md}):

\begin{itemize}[leftmargin=*]
  \item Renderizar A4 a 500 DPI genera imágenes de $\sim$\,4000$\times$5800 px ($\approx$\,70 MB
        en RAM por página). Procesar página por página.
  \item PDFs con texto vectorizado obligan al uso de OCR (no admiten extracción directa).
  \item PDFs con texto embebido pueden procesarse más rápido con \code{PDFTextStripper} sin
        Tesseract; el módulo actual asume el caso vectorizado.
\end{itemize}

% =====================================================================
\section{Frontend (Vite + Vanilla JS)}
% =====================================================================

\subsection{Decisiones de diseño}

El frontend se implementó en JavaScript vanilla sobre Vite, sin framework SPA. Esta decisión
prioriza:

\begin{itemize}[leftmargin=*]
  \item \textbf{Cero dependencias de runtime} más allá de Chart.js.
  \item \textbf{Curva de mantenimiento baja} para un proyecto académico.
  \item \textbf{Tiempo de carga mínimo} y \emph{bundle} reducido.
\end{itemize}

\subsection{Cliente HTTP (\code{api.js})}

\begin{itemize}[leftmargin=*]
  \item Encapsula \code{fetch} con manejo de \code{Authorization: Bearer <token>}.
  \item Persiste el access token y el refresh token en \code{localStorage}.
  \item \textbf{Auto-desempaqueta} la propiedad \code{.data} del envoltorio
        \code{ApiResponse<T>} para todas las respuestas de dominio. Las respuestas de
        \code{/api/auth/*} no están envueltas y se devuelven tal cual.
  \item Detecta 401 y dispara el flujo de \emph{refresh}.
\end{itemize}

\subsection{Páginas y permisos por rol}

\begin{table}[h]\centering\small\renewcommand{\arraystretch}{1.15}
\begin{tabular}{@{}p{4.6cm} p{2.0cm} p{6.5cm}@{}}
\toprule
\textbf{Página} & \textbf{Rol} & \textbf{Función} \\
\midrule
\code{aniosLectivos.js}      & ADMIN/COORD       & Gestión de años lectivos y trimestres \\
\code{secciones.js}          & ADMIN/COORD       & Gestión de secciones y asignaciones \\
\code{estudiantes.js}        & Todos             & Listado y CRUD de estudiantes \\
\code{importarPiad.js}       & DOCENTE+          & Carga de PDF PIAD para extracción OCR \\
\code{evaluaciones.js}       & DOCENTE+          & Evaluaciones 4-puntos (legado) \\
\code{evaluacionesSaber.js}  & DOCENTE+          & Wizard de evaluación por saber \\
\code{matriculas.js}         & DOCENTE+          & Matrícula de estudiantes en secciones \\
\code{alertas.js}            & DOCENTE+          & Listado de alertas tempranas (legado) \\
\code{alertasTempranas.js}   & DOCENTE+          & Alertas temáticas con botón de impresión \\
\code{visualizaciones.js}    & DOCENTE+          & Gráficas Chart.js de distribución \\
\code{reportes.js}           & DOCENTE+          & Reportes por estudiante \\
\code{sesion.js}             & Todos             & Sección de sesión en el menú lateral \\
\bottomrule
\end{tabular}
\caption{Páginas del frontend y rol mínimo requerido (ADMIN/COORD = ADMIN o COORDINADOR).}
\end{table}

\subsection{Mejoras de UX implementadas}

\begin{itemize}[leftmargin=*]
  \item \textbf{Confirmación al cerrar wizard} con datos sin guardar
        (\code{evaluacionesSaber.js}).
  \item \textbf{Botón “Aplicar nivel a todos los ejes”} en el wizard de saberes para
        evaluación rápida cuando el desempeño es uniforme.
  \item \textbf{Botón de impresión} en Alertas Tempranas que genera una vista lista para PDF.
  \item \textbf{Menú lateral filtrado por rol:} los docentes solo ven Estudiantes, Importar
        PIAD, Evaluaciones por Saber, Alertas Tempranas, Visualizaciones y Reportes.
        Años Lectivos y Secciones aparecen únicamente para ADMIN/COORDINADOR.
  \item \textbf{Inicio contextual por rol:} los docentes entran directamente a
        \emph{Evaluaciones por Saber}; los administradores y coordinadores aterrizan en
        \emph{Años Lectivos}.
\end{itemize}

% =====================================================================
\section{Infraestructura de pruebas funcionales}
% =====================================================================

\subsection{Descripción}

Las pruebas funcionales validan el sistema completo de extremo a extremo: desde la petición
HTTP hasta la base de datos. \textbf{No usan \emph{mocks}}; ejecutan contra una instancia real
de PostgreSQL levantada con \textbf{Testcontainers}.

\subsection{Arquitectura de las pruebas}

\begin{lstlisting}[style=bash,caption={Estructura de las clases de prueba funcional}]
IntegrationTestBase (abstracta)
   |-- AuthFuncionalTest             (10 tests)
   |-- ControlAccesoFuncionalTest    (10 tests)
   |-- EvaluacionSaberFuncionalTest  ( 8 tests)
   `-- AlertaTematicaFuncionalTest   ( 5 tests)
                                    -----
                          Total:    33 tests
\end{lstlisting}

Adicionalmente existen pruebas unitarias clásicas a nivel de servicio (sin Docker):
\code{EstudianteServiceImplTest} y \code{PeriodoServiceImplTest}.

\subsection{Características técnicas}

\begin{itemize}[leftmargin=*]
  \item \textbf{Base de datos real:} \code{PostgreSQLContainer} con la imagen
        \code{postgres:15-alpine}. Se aplica el patrón \emph{singleton container}: el
        contenedor arranca una vez para toda la sesión JVM y se detiene al finalizar
        (\emph{shutdown hook} automático de Testcontainers).
  \item \textbf{Flyway automático:} Spring ejecuta las migraciones V1--V11 al levantar el
        contexto, garantizando datos semilla consistentes en cada corrida.
  \item \textbf{Sin \code{@Transactional}:} con \code{RANDOM\_PORT} el servidor corre en un
        hilo separado y los \emph{commits} son reales. El aislamiento se logra con
        \code{@AfterEach} que llama a \code{limpiarEvaluacionesSaber(long... estudianteIds)}
        a través de \code{JdbcTemplate}.
  \item \textbf{\code{RestTemplate} con \code{NoOpResponseErrorHandler}:} no lanza excepción
        en respuestas 4xx/5xx, permitiendo verificar el código HTTP directamente.
  \item \textbf{AssertJ:} aserciones fluidas sobre el estado HTTP y el cuerpo JSON.
\end{itemize}

\subsection{Resumen de resultados}

\begin{table}[h]\centering\small\renewcommand{\arraystretch}{1.2}
\begin{tabular}{@{}l c c c c@{}}
\toprule
\textbf{Clase de prueba} & \textbf{Tests} & \textbf{Pass} & \textbf{Fail} & \textbf{Error} \\
\midrule
\code{AuthFuncionalTest}            & 10  & 10  & 0 & 0 \\
\code{ControlAccesoFuncionalTest}   & 10  & 10  & 0 & 0 \\
\code{EvaluacionSaberFuncionalTest} &  8  &  8  & 0 & 0 \\
\code{AlertaTematicaFuncionalTest}  &  5  &  5  & 0 & 0 \\
\midrule
\textbf{TOTAL}                      & \textbf{33} & \textbf{33} & \textbf{0} & \textbf{0} \\
\bottomrule
\end{tabular}
\end{table}

\subsection{Salida observada del \emph{build}}

\begin{ataraExito}
\begin{lstlisting}[style=bash,basicstyle=\ttfamily\scriptsize]
$ ./mvnw test -Dtest="*Funcional*"

[INFO] Tests run:  5  Failures: 0  Errors: 0  Skipped: 0  (AlertaTematica)
[INFO] Tests run: 10  Failures: 0  Errors: 0  Skipped: 0  (Autenticacion)
[INFO] Tests run: 10  Failures: 0  Errors: 0  Skipped: 0  (ControlAcceso)
[INFO] Tests run:  8  Failures: 0  Errors: 0  Skipped: 0  (EvaluacionSaber)
[INFO] Tests run: 33  Failures: 0  Errors: 0  Skipped: 0

[INFO] BUILD SUCCESS
\end{lstlisting}
\end{ataraExito}

% =====================================================================
\section{Configuración y despliegue}
% =====================================================================

\subsection{Comandos de desarrollo local}

\begin{lstlisting}[style=bash,caption={Comandos esenciales}]
# Levantar PostgreSQL (requerido antes de la app)
docker-compose up -d

# Ejecutar la aplicacion (puerto 8081)
./mvnw spring-boot:run

# Compilar
./mvnw clean package

# Ejecutar todas las pruebas (Docker requerido)
./mvnw test

# Una sola clase de prueba
./mvnw test -Dtest=NombreClase

# Un solo metodo
./mvnw test -Dtest=NombreClase#nombreMetodo

# Frontend de desarrollo (puerto 3000)
cd temp-frontend && npm install && npm run dev
\end{lstlisting}

\subsection{Variables de entorno}

\begin{table}[h]\centering\small\renewcommand{\arraystretch}{1.15}
\begin{tabular}{@{}p{4.4cm} p{4.0cm} p{6.5cm}@{}}
\toprule
\textbf{Variable} & \textbf{Default (dev)} & \textbf{Descripción} \\
\midrule
\code{DB\_URL}                       & \code{jdbc:postgresql://localhost:5433/atara\_db} & URL de la base de datos \\
\code{DB\_USERNAME}                  & \code{atara\_user}              & Usuario de la BD \\
\code{DB\_PASSWORD}                  & \code{atara\_pass\_123}         & Contraseña de la BD \\
\code{JWT\_SECRET}                   & cadena de desarrollo            & Clave JWT — \textbf{cambiar} en producción ($\geq$ 64 chars) \\
\code{JWT\_EXPIRATION\_MS}           & \code{86400000} (1 día)          & Duración del access token \\
\code{JWT\_REFRESH\_EXPIRATION\_DAYS}& \code{7}                         & Duración del refresh token \\
\code{CORS\_ALLOWED\_ORIGINS}        & \code{*}                         & Origen(es) permitidos del frontend \\
\code{SHOW\_SQL}                     & \code{true}                     & Logs de SQL — \code{false} en producción \\
\code{SPRING\_PROFILES\_ACTIVE}      & ---                              & Establecer \code{prod} para activar \code{application-prod.properties} \\
\code{TESSDATA\_PREFIX}              & ---                              & Ruta alterna a \code{tessdata} (PIAD/OCR) \\
\bottomrule
\end{tabular}
\caption{Variables de entorno admitidas con \emph{fallback} para desarrollo local.}
\end{table}

\subsection{Perfil de producción}

El archivo \code{application-prod.properties} se activa con
\code{SPRING\_PROFILES\_ACTIVE=prod}. Sus diferencias respecto a desarrollo:

\begin{itemize}[leftmargin=*]
  \item \code{spring.jpa.show-sql=false} y \code{format\_sql=false} (sin logs de SQL).
  \item Todas las propiedades sensibles se obtienen exclusivamente de variables de entorno.
  \item Se exige \code{JWT\_SECRET} de al menos 64 caracteres aleatorios.
  \item \code{CORS\_ALLOWED\_ORIGINS} debe apuntar al dominio real (por ejemplo
        \code{https://atara.mep.go.cr}), nunca \code{*}.
\end{itemize}

\subsection{Despliegue con Docker (BD)}

\begin{lstlisting}[style=bash,caption={\code{docker-compose.yml} de desarrollo}]
services:
  postgres:
    image: postgres:17
    container_name: atara-postgres
    restart: unless-stopped
    environment:
      POSTGRES_DB: atara_db
      POSTGRES_USER: atara_user
      POSTGRES_PASSWORD: atara_pass_123
    ports:
      - "5433:5432"
    volumes:
      - atara_postgres_data:/var/lib/postgresql/data
    healthcheck:
      test: ["CMD-SHELL", "pg_isready -U atara_user -d atara_db"]
      interval: 10s
      timeout: 5s
      retries: 5
volumes:
  atara_postgres_data:
\end{lstlisting}

% =====================================================================
\section{Mapa completo de endpoints}
% =====================================================================

\renewcommand{\arraystretch}{1.15}\small
\begin{longtable}{@{}p{5.5cm} p{1.6cm} >{\raggedright\arraybackslash}p{8.5cm}@{}}
\toprule
\textbf{Controlador} & \textbf{Método} & \textbf{Ruta} \\
\midrule
\endfirsthead

\toprule
\textbf{Controlador} & \textbf{Método} & \textbf{Ruta} \\
\midrule
\endhead

\bottomrule
\multicolumn{3}{r}{\itshape (continúa en la siguiente página)}
\endfoot

\bottomrule
\endlastfoot

\code{AuthController}             & POST    & \ruta{/api/auth/login} \\
                                  & POST    & \ruta{/api/auth/refresh} \\
                                  & POST    & \ruta{/api/auth/logout} \\
                                  & GET     & \ruta{/api/auth/me} \\
\midrule
\code{EstudianteController}       & POST    & \ruta{/api/estudiantes} \\
                                  & GET     & \ruta{/api/estudiantes} \\
                                  & GET     & \ruta{/api/estudiantes/{id}} \\
                                  & PUT     & \ruta{/api/estudiantes/{id}} \\
\midrule
\code{MatriculaController}        & POST    & \ruta{/api/matriculas} \\
                                  & GET     & \ruta{/api/matriculas/estudiante/{estudianteId}} \\
                                  & GET     & \ruta{/api/matriculas/seccion/{seccionId}} \\
\midrule
\code{AnioLectivoController}      & POST    & \ruta{/api/anios-lectivos} \\
                                  & GET     & \ruta{/api/anios-lectivos} \\
                                  & GET     & \ruta{/api/anios-lectivos/activo} \\
                                  & GET     & \ruta{/api/anios-lectivos/{id}} \\
                                  & PUT     & \ruta{/api/anios-lectivos/{id}} \\
                                  & PUT     & \ruta{/api/anios-lectivos/{id}/activar} \\
                                  & DELETE  & \ruta{/api/anios-lectivos/{id}} \\
\midrule
\code{PeriodoController}          & POST    & \ruta{/api/periodos} \\
                                  & GET     & \ruta{/api/periodos} \\
                                  & GET     & \ruta{/api/periodos/activo} \\
                                  & PUT     & \ruta{/api/periodos/{id}} \\
                                  & PUT     & \ruta{/api/periodos/{id}/activar} \\
                                  & DELETE  & \ruta{/api/periodos/{id}} \\
\midrule
\code{SeccionController}          & GET     & \ruta{/api/secciones} \\
                                  & GET     & \ruta{/api/secciones/docente/{docenteId}} \\
                                  & POST    & \ruta{/api/secciones} \\
                                  & PUT     & \ruta{/api/secciones/{id}} \\
                                  & DELETE  & \ruta{/api/secciones/{id}} \\
                                  & GET     & \ruta{/api/secciones/catalogos/niveles} \\
                                  & GET     & \ruta{/api/secciones/catalogos/centros} \\
                                  & GET     & \ruta{/api/secciones/catalogos/docentes} \\
\midrule
\code{CatalogoSaberController}    & GET     & \ruta{/api/catalogos/saberes/tipos} \\
                                  & GET     & \ruta{/api/catalogos/saberes/materias} \\
                                  & GET     & \ruta{/api/catalogos/saberes/ejes} \\
                                  & GET     & \ruta{/api/catalogos/saberes/niveles-desempeno} \\
\midrule
\code{EvaluacionSaberController}  & POST    & \ruta{/api/evaluaciones-saber} \\
                                  & PUT     & \ruta{/api/evaluaciones-saber/{id}} \\
                                  & GET     & \ruta{/api/evaluaciones-saber/{id}} \\
                                  & GET     & \ruta{/api/evaluaciones-saber/estudiante/{eId}/periodo/{pId}} \\
                                  & GET     & \ruta{/api/evaluaciones-saber/seccion/{sId}/periodo/{pId}} \\
                                  & GET     & \ruta{/api/evaluaciones-saber/promedios/estudiante/{eId}/periodo/{pId}} \\
                                  & GET     & \ruta{/api/evaluaciones-saber/promedios/seccion/{sId}/periodo/{pId}} \\
\midrule
\code{AlertaTematicaController}   & POST    & \ruta{/api/alertas-tematicas/generar/estudiante/{eId}/periodo/{pId}} \\
                                  & POST    & \ruta{/api/alertas-tematicas/generar/seccion/{sId}/periodo/{pId}} \\
                                  & GET     & \ruta{/api/alertas-tematicas/estudiante/{eId}/periodo/{pId}} \\
                                  & GET     & \ruta{/api/alertas-tematicas/seccion/{sId}/periodo/{pId}} \\
\midrule
\code{EvaluacionController}       & POST    & \ruta{/api/evaluaciones} \\
                                  & POST    & \ruta{/api/evaluaciones/{id}/detalles} \\
                                  & GET     & \ruta{/api/evaluaciones/{id}} \\
                                  & GET     & \ruta{/api/evaluaciones/estudiante/{estudianteId}} \\
                                  & GET     & \ruta{/api/evaluaciones/estudiante/{estudianteId}/periodo/{periodoId}} \\
                                  & GET     & \ruta{/api/evaluaciones/periodo/{periodoId}} \\
\midrule
\code{AlertaController}           & GET     & \ruta{/api/alertas/estudiante/{studentId}} \\
                                  & GET     & \ruta{/api/alertas/seccion/{sectionId}} \\
\midrule
\code{ReporteController}          & GET     & \ruta{/api/reportes/estudiante/{studentId}} \\
\midrule
\code{VisualizacionController}    & GET     & \ruta{/api/visualizaciones/seccion/{sectionId}/distribucion} \\
\midrule
\code{PiadController}             & POST    & \ruta{/api/piad/extraer} \\

\end{longtable}

% =====================================================================
\section{Convenciones de desarrollo}
% =====================================================================

\subsection{Fronteras entre capas}

\textbf{Los controladores deben:}

\begin{itemize}[leftmargin=*]
  \item Recibir requests, delegar al servicio y retornar respuesta.
  \item Declarar \code{@Valid} en el cuerpo de la petición.
  \item Retornar \code{ResponseEntity<ApiResponse<T>>}.
\end{itemize}

\textbf{Los controladores nunca deben:}

\begin{itemize}[leftmargin=*]
  \item Contener lógica de negocio (\code{if/else}, cálculos, transformaciones).
  \item Llamar repositorios directamente.
  \item Retornar entidades JPA o tipos internos del modelo.
\end{itemize}

\subsection{DTOs obligatorios}

\begin{itemize}[leftmargin=*]
  \item \textbf{Nunca} retornar ni aceptar una entidad JPA en un controlador.
  \item Los DTO se ubican en \code{dto/<dominio>/} (\code{EstudianteRequestDto.java},
        \code{EstudianteResponseDto.java}, etc.).
  \item Mapeo entidad$\leftrightarrow$DTO en la capa de servicio.
  \item DTOs con Lombok: \code{@Data}, \code{@Builder}, \code{@NoArgsConstructor},
        \code{@AllArgsConstructor}.
\end{itemize}

\subsection{Validación}

\begin{itemize}[leftmargin=*]
  \item Anotaciones de Jakarta Bean Validation: \code{@NotNull}, \code{@NotBlank}, \code{@Size},
        \code{@Min}, \code{@Max}, \code{@Email}, \code{@Pattern}, etc.
  \item \code{@Valid} en el parámetro del método controlador para activar la validación.
  \item No se permite validación manual cuando una anotación cubre el caso.
  \item \code{GlobalExceptionHandler} maneja \code{MethodArgumentNotValidException} y
        retorna errores estructurados a nivel de campo (\code{campo: mensaje}).
\end{itemize}

\subsection{Estructura de respuesta API}

Todos los endpoints retornan \code{ResponseEntity<ApiResponse<T>>} con la siguiente forma:

\begin{lstlisting}[style=json,caption={Forma del envoltorio \code{ApiResponse} (éxito y error)}]
// Éxito
{ "success": true,  "message": "Descripcion breve", "data": { ... } }

// Error
{ "success": false, "message": "Descripcion del error", "data": null }
\end{lstlisting}

\subsection{Manejo centralizado de errores}

\begin{itemize}[leftmargin=*]
  \item Toda excepción pasa por \code{GlobalExceptionHandler}: prohibidos los \code{try/catch}
        en controladores.
  \item Mapeo:
  \begin{itemize}
    \item \code{NoSuchElementException} $\to$ 404 Not Found.
    \item \code{IllegalArgumentException} $\to$ 400 Bad Request.
    \item \code{AccesoDenegadoException} $\to$ 403 Forbidden.
    \item \code{TokenRefreshException} $\to$ 401 Unauthorized.
    \item \code{MethodArgumentNotValidException} $\to$ 400 con detalle por campo.
    \item Otras $\to$ 500 Internal Server Error (sin exponer \emph{stacktrace}).
  \end{itemize}
\end{itemize}

\subsection{Convenciones de nomenclatura}

\begin{table}[h]\centering\small\renewcommand{\arraystretch}{1.2}
\begin{tabular}{@{}p{4.6cm} p{4.0cm} p{4.5cm}@{}}
\toprule
\textbf{Elemento} & \textbf{Convención} & \textbf{Ejemplo} \\
\midrule
Clases                  & PascalCase           & \code{EstudianteServiceImpl} \\
Métodos / variables     & camelCase            & \code{obtenerEstudiantePorId} \\
Constantes              & UPPER\_SNAKE\_CASE   & \code{MAX\_INTENTOS\_LOGIN} \\
Columnas BD             & snake\_case          & \code{fecha\_nacimiento} \\
DTO                     & sufijo \code{RequestDto} / \code{ResponseDto} & \code{MatriculaRequestDto} \\
Interfaz de servicio    & sin sufijo           & \code{EstudianteService} \\
Implementación servicio & sufijo \code{Impl}   & \code{EstudianteServiceImpl} \\
Rutas REST              & kebab-case, plural   & \code{/api/anios-lectivos} \\
\bottomrule
\end{tabular}
\caption{Convenciones de nomenclatura adoptadas en el proyecto.}
\end{table}

\subsection{Reglas de seguridad para desarrolladores}

\begin{ataraAlerta}
\begin{itemize}[leftmargin=*]
  \item Nunca registrar contraseñas, tokens ni PII en logs (cédulas, nombres con calificaciones).
  \item Los tokens JWT \emph{solo} se validan en \code{JwtAuthenticationFilter}.
  \item No exponer IDs internos ni mensajes técnicos en los errores cara al cliente.
  \item Las contraseñas se almacenan con \code{BCryptPasswordEncoder} (\emph{strength} 12).
  \item Sanitizar toda cadena del usuario antes de usarla en JPQL/SQL nativo; preferir
        \emph{method names} de Spring Data o \code{@Query} con parámetros nombrados.
  \item Los servicios de dominio deben recibir el \code{ContextoUsuario} como parámetro
        explícito; está \textbf{prohibido} llamar a \code{SecurityContextHolder} directamente
        fuera de \code{ContextoUsuarioService}.
\end{itemize}
\end{ataraAlerta}

% =====================================================================
\section{Glosario}
% =====================================================================

\begin{description}[leftmargin=2.0cm,style=nextline,labelwidth=*,font=\bfseries\color{ataraAzul}]
  \item[ATARA] Sistema de Alerta Temprana y Atención al Rendimiento Académico.
  \item[Eje temático] Subdivisión pedagógica de un \emph{tipo de saber} dentro de una materia.
                      Cada materia tiene 21 ejes (7 por tipo de saber).
  \item[Tipo de saber] Categoría pedagógica del aprendizaje: Conceptual, Procedimental o
                       Actitudinal.
  \item[Materia] Asignatura curricular: Español, Matemáticas, Ciencias o Estudios Sociales.
  \item[Periodo] Trimestre lectivo (I, II, III) dentro de un \emph{año lectivo}.
  \item[ContextoUsuario] \emph{Record} inmutable que encapsula la identidad y el alcance
                         del usuario autenticado.
  \item[Alerta temática] Registro generado automáticamente cuando el promedio de un eje
                         temático cae bajo el umbral (ALTA $\leq 2{,}0$, MEDIA $2{,}1$--$3{,}0$).
  \item[PIAD] Plan Individual de Adecuación: documento oficial del MEP que lista a los
              estudiantes con adecuaciones curriculares.
  \item[Flyway] Herramienta de migraciones de esquema usada para la gestión de la base de
                datos.
  \item[Testcontainers] Librería que arranca contenedores Docker como dependencia de pruebas
                        JUnit.
  \item[ApiResponse] Envoltorio uniforme \code{\{success, message, data\}} para todas las
                     respuestas del API.
\end{description}

% =====================================================================
\section{Conclusiones}
% =====================================================================

\begin{itemize}[leftmargin=*]
  \item ATARA implementa una API REST estructurada en \textbf{14 controladores} y
        \textbf{7 módulos funcionales principales}, siguiendo una arquitectura por capas
        con separación estricta de responsabilidades y un cliente frontend independiente.
  \item El sistema de \textbf{control de acceso por rol}, basado en el \code{ContextoUsuario},
        garantiza que cada docente opera exclusivamente sobre sus secciones y materias
        asignadas, sin requerir lógica condicional dispersa en cada endpoint.
  \item El módulo de \textbf{alertas temáticas} automatiza la detección de riesgo académico:
        un docente registra evaluaciones por saber y el sistema identifica los ejes que
        requieren atención inmediata (ALTA), seguimiento (MEDIA) o ninguna acción.
  \item La suite de \textbf{33 pruebas funcionales} se ejecuta contra una base de datos
        PostgreSQL real (Testcontainers) con datos semilla completos, validando flujos de
        negocio extremo a extremo con un \textbf{100\,\% de éxito}.
  \item La \textbf{migración incremental con Flyway} (V1--V11) garantiza la reproducibilidad
        del esquema en cualquier entorno (desarrollo, pruebas, producción) y mantiene un
        historial auditable de cambios.
  \item La \textbf{externalización de credenciales} a variables de entorno y el perfil
        \code{prod} permiten un despliegue seguro sin modificar el código fuente.
  \item El módulo \textbf{PIAD/OCR} demuestra integración avanzada (PDFBox + Tess4j) para
        extraer datos de documentos vectorizados; la persistencia automática de la lista
        extraída queda como trabajo futuro (V12).
\end{itemize}

\vspace{1em}

\begin{ataracaja}[title=Trabajo futuro recomendado]
\begin{itemize}[leftmargin=*]
  \item \textbf{V12 — Persistencia PIAD:} guardar automáticamente los estudiantes extraídos
        por OCR en \code{POST /api/piad/extraer}, reutilizando la idempotencia por cédula.
  \item Ampliar la cobertura de pruebas funcionales a los módulos de \textbf{Reportes} y
        \textbf{Visualizaciones}.
  \item Añadir excepciones específicas \code{RecursoNoEncontradoException} y
        \code{ReglaDeNegocioException} para reemplazar el uso genérico de
        \code{NoSuchElementException} y \code{IllegalArgumentException}.
  \item Implementar paginación (\code{Pageable}) en los listados largos
        (\code{GET /api/estudiantes}, \code{GET /api/secciones}).
  \item Migrar el frontend a un \emph{bundle} estable de producción servido por el mismo
        proceso Spring (o desplegado como artefacto estático en CDN).
\end{itemize}
\end{ataracaja}

\end{document}
